\exnumber{01}
\extitle{UI II - Ex 2 - Material Design}
\exfiles{ex01/fichier Figma, nomme Name\_FirstName\_DeliverableName\_Date\_VersionNumber}
\exauthorize{Ceux nommés dans l'exercice}
\makeheaderfiles

\textbf{Instructions :}

Téléchargez cette image \url{https://storage.googleapis.com/gd-wagtail-prod-assets/original_images/MDA2018_inline_02.jpg} qui est un aperçu de 3 écrans de Lyft, et détectez l'ordre dans lequel les calques ont été posés.

\textbf{Livrables :}

\begin{itemize}
\item Un fichier Figma avec l'image
\item Sur chaque écran, indiquer l'ordre des calques avec des numéros. 1 étant l'image de fond, 2 l'élément juste au-dessus, etc.
\item Déposez le fichier Figma sur Github, nommé ainsi : « Name\_FirstName\_DeliverableName\_Date\_VersionNumber ». Exemple : « Doe\_John\_MidFiPrototype\_05242024\_V1 ».
\item Ajoutez un titre dans le fichier.
\item Soignez l'aspect global du fichier. Il doit être clair, simple et lisible. Vous pouvez regarder canva.com, sans surcharger le dessin de détails.
\end{itemize}

\textbf{Pistes :}

\begin{itemize}
\item 2 chiffres peuvent être intervertis tant que l'ordre des calques est respecté
\end{itemize}

\textbf{Ressources :}

\begin{itemize}
\item Material Design Awards 2019 \url{https://design.google/library/material-design-awards-2019/}
\item Qu'est-ce que Material Design - Wikipedia \url{https://en.wikipedia.org/wiki/Material_Design}
\end{itemize}

